% Part: first-order-logic
% Chapter: model-theory
% Section: overspill

\documentclass[../../../include/open-logic-section]{subfiles}

\begin{document}

\olfileid{mod}{bas}{ove}
\olsection{सांत मर्यादेपलीकडील प्रसरण}

\begin{thm}
\ollabel{overspill} वाक्यांचा एखादा संच $\Gamma$ याला हवे तितके मोठे
सांत प्रतिमान असतील, तर त्याला एक अनंत प्रतिमान असते.
\end{thm}

\begin{proof}
$c_0$, $c_1$, \dots ही गणनीय इतकी नवीन स्थिरे जोडून $\Gamma$ च्या भाषेचा
विस्तार करा आणि $\Gamma \cup \{c_i \neq c_j : i \neq j\}$ हा संच विचारात
घ्या. $\Gamma$ ला हवे तितके मोठे सांत प्रतिमान आहेत, याचा अर्थ प्रत्येक
$m >0$ साठी असा $n\ge m$ असतो की $\Gamma$ ला आकारमान~$n$ चे प्रतिमान
असते. यावरून $\Gamma \cup \{c_i \neq c_j : i \neq j\}$ सांततः
पूर्ततायोग्य आहे. संहततेनुसार $\Gamma \cup \{c_i \neq c_j : i \neq j\}$
याला प्रतिमान $\Struct M$ असते; त्याचा प्रांत अनंत असलाच पाहिजे, कारण ते
$c_i \neq c_j$ या सर्व असमानतांची पूर्ति करते.
\end{proof}

\begin{prop}
\ollabel{inf-not-fo}
कोणत्याही प्रथम-क्रम भाषेत असे कोणतेही वाक्य $!A$ नाही की ते
!!a{structure}~$\Struct M$ मध्ये तेव्हा आणि केवळ तेव्हाच सत्य असते, जेव्हा
त्या !!{structure} चा प्रांत $\Domain{M}$ अनंत असतो.
\end{prop}

\begin{proof}
असे एखादे $!A$ असते, तर त्याचे नकरण $\lnot !A$ सर्व आणि फक्त सांत
!!{structure}s मध्ये सत्य असते. म्हणून त्याला हवे तितके मोठे सांत प्रतिमान
असतात, पण अनंत प्रतिमान नसते; हे \olref{overspill} शी विसंगत आहे.
\end{proof}

\end{document}
